\chapter{规划、委托与多 Agent}
\label{chap:multi-agent}

\begin{takeaway}[学习目标]
理解显式计划的价值与边界；判断何时需要子 Agent；设计任务委托合同、并行合并、隔离和失败传播策略，并用数据证明多 Agent 是否值得。
\end{takeaway}

\section{计划是可修改的工作假设}

规划的价值不在于让模型提前写一篇长推理，而在于把任务拆成可观察状态：下一步是什么、依赖什么、怎样验收、失败后如何调整。环境发生变化时，计划必须能重写。

简单任务可以边走边决定；长任务适合维护结构化计划。一个步骤至少包含：ID、目标、依赖、执行者、输入、预期产物、状态和验证方式。

\begin{codeblock}
{
  "id": "research-pricing",
  "goal": "核验 5 个产品当前公开价格",
  "depends_on": ["select-products"],
  "output": "pricing.json",
  "acceptance": "每个价格含官方来源和访问日期",
  "budget": {"minutes": 15, "requests": 30}
}
\end{codeblock}

\section{何时重新规划}

触发条件应由系统显式定义：关键假设被工具否定；依赖资源不可用；剩余预算不足；连续失败；用户改变目标；新证据使后续步骤失效。不要每一步都全量重写计划，否则容易产生规划抖动和额外成本。

\section{多 Agent 的真实收益}

多 Agent 只有在下列至少一项成立时才可能值得：

\begin{itemize}
  \item 子任务可以真正并行，显著降低总时间。
  \item 不同子任务需要互相冲突的上下文或工具权限，隔离能提升表现。
  \item 单个上下文装不下完整任务，委托能压缩边界。
  \item 需要独立证据或对抗审查，分离执行者与评审者可减少同源偏差。
\end{itemize}

“给三个角色起名字”不是收益。多 Agent 会新增调用成本、通信损失、状态一致性、重复工作、合并冲突和调试难度。

\section{委托合同}

主 Agent 给 worker 的消息应像一份小任务合同，而不是“帮我研究一下”：

\begin{enumerate}
  \item 唯一目标与不在范围内的事项；
  \item 已知输入、来源范围和权限；
  \item 产物 schema 与保存位置；
  \item 可验证的完成条件；
  \item 时间、token、工具和调用次数预算；
  \item 失败、冲突和需要升级的条件。
\end{enumerate}

只传完成子任务所需的上下文。把主任务全部历史复制给每个 worker 会增加成本，也会把无关指令和敏感数据扩散到更多执行单元。

\section{并行任务的合并}

合并器不能只把输出拼接起来。它应执行 schema 校验、去重、来源核验、冲突识别、覆盖度检查和最终排序。若两个 worker 对同一事实冲突，应保留各自证据并触发核验任务，不能让模型凭语气选一个。

\begin{table}[htbp]
\centering
\begin{tabularx}{\textwidth}{>{\bfseries}p{0.2\textwidth} X X}
\toprule
问题 & 典型原因 & 控制方法 \\
\midrule
重复劳动 & 边界不清、共享发现不可见 & 分区键、共享索引、任务锁 \\
输出冲突 & schema 或口径不同 & 统一合同、证据优先的合并器 \\
级联失败 & 主任务依赖单个 worker & 超时、替代 worker、部分结果 \\
成本失控 & 递归委托或无预算 & 委托深度、全局预算、配额 \\
权限扩散 & worker 继承全部工具 & 最小权限、临时凭证、沙箱 \\
\bottomrule
\end{tabularx}
\end{table}

\section{共享状态与隔离状态}

共享状态适合任务队列、已确认事实和 artifact 索引；隔离状态适合 worker 草稿、私有上下文和临时凭证。共享写入要有版本或事务机制，避免最后写入者覆盖别人结果。

推荐使用“追加事件 + 派生视图”：worker 追加发现事件，编排器生成当前状态视图。事件含来源、作者、时间和关联步骤，便于审计与回放。

\section{评审 Agent 也会犯错}

让另一个模型评分可以扩展评测，但不代表客观。评审者可能偏爱更长答案、受参考答案措辞影响，或与生成模型共享盲点。降低风险的方法包括：明确 rubric、隐藏无关身份、随机答案顺序、使用程序化断言、抽样人工校准，以及记录评审理由。

\section{用实验决定是否拆分}

对同一测试集比较单 Agent 与多 Agent：任务完成率、端到端时间、总 token、工具调用数、人工修订时间、故障可定位性。多 Agent 如果只提升了“看起来很复杂”，就应该删除。

\begin{practice}[并行竞品研究]
实现一个编排器，把五个竞品按产品分给 worker。每个 worker 只返回统一 JSON：事实、来源、访问日期、未知项。合并器必须发现字段缺失和价格冲突。与单 Agent 串行版本比较总时间和事实准确率。
\end{practice}

\begin{checkpoint}
\begin{itemize}
  \item \checkbox 计划步骤有依赖、预算、产物和验收方式。
  \item \checkbox 每个委托都有边界，worker 只获得最小上下文和权限。
  \item \checkbox 合并器会校验、去重并暴露证据冲突。
  \item \checkbox 我用基准数据证明了多 Agent 的净收益。
\end{itemize}
\end{checkpoint}

